% CAISc 2026 submission — Physics Auditor paper
% Uses official caisc_2026 style file (submission mode = anonymous + line numbers)

\documentclass{article}

\usepackage{caisc_2026}      % submission mode: anonymous + line numbers

\usepackage[utf8]{inputenc}
\usepackage[T1]{fontenc}
\usepackage{hyperref}
\usepackage{url}
\usepackage{booktabs}
\usepackage{amsfonts}
\usepackage{amsmath}
\usepackage{amssymb}
\usepackage{nicefrac}
\usepackage{microtype}
\usepackage{xcolor}
\usepackage{graphicx}

\title{AI Scientists Cannot Self-Diagnose Their Own Physics Violations:\\
A Physics Auditor Agent for Conservation Law Verification\\
in AI-Generated PINN Experiments}

% Author block: auto-anonymized by caisc_2026 in submission mode.
% List human authors only per CAISc authorship policy.
% AI agents are documented in the AI Involvement Checklist below.
\author{}

\begin{document}

\maketitle

\begin{abstract}
We present the first systematic study of AI scientists critiquing
AI-generated science: a Physics Auditor agent that automatically detects
conservation law violations in PINN experiments produced by end-to-end AI
research systems, revealing that state-of-the-art AI scientists generate
physically inconsistent results at high rates that neither their built-in AI
reviewers nor prompted LLM baselines reliably detect.
End-to-end AI scientist systems such as the AI Scientist
\citep{Lu2024AIScientist} and Agent Laboratory \citep{Wang2024AgentLab}
close the loop from hypothesis generation through code execution to write-up
and peer review, yet none implement domain-specific physical consistency
checks on their outputs.
We evaluate $N=40$ synthetic AI-scientist-style PINN outputs spanning five
violation types drawn from the published PINN failure taxonomy.
The Physics Auditor achieves 100\% recall (95\% CI: [0.893, 1.000]) on all
injected violations; a prompted LLM baseline reviewer achieves 46.9\% recall
(95\% CI: [0.309, 0.635]) --- a gap of 53.1 percentage points
(McNemar's test: $\chi^2 = 15.06$, $p = 0.0001$).
Two violation types --- mass non-conservation and PDE residual error ---
are structural blind spots for LLM reviewers, each yielding 0\% detection rate.
All violation scores are auto-computable from the released dataset;
any reviewer can re-run the Physics Auditor on the 40 test outputs and
reproduce Table~\ref{tab:main} in under two minutes on a standard laptop.
\end{abstract}

%%%%%%%%%%%%%%%%%%%%%%%%%%%%%%%%%%%%%%%%%%%%%%%%%%%%%%%%%%%%
\section{Introduction}
%%%%%%%%%%%%%%%%%%%%%%%%%%%%%%%%%%%%%%%%%%%%%%%%%%%%%%%%%%%%

End-to-end AI scientist systems have reached the point where they generate,
execute, and peer-review physics experiments with minimal human intervention
--- yet none of them verify whether the results obey physics.
The AI Scientist \citep{Lu2024AIScientist} and its successor AI Scientist-v2
\citep{Lu2025AIScientistV2} automate the full scientific loop ---
hypothesis, code, experiment, write-up, and peer review --- for machine
learning research including PINN-based PDE tasks.
Agent Laboratory \citep{Wang2024AgentLab} and ResearchAgent
\citep{Baek2024ResearchAgent} extend this paradigm to broader scientific domains.
These systems are generating results with minimal human intervention.

The problem is structural: none of these systems check whether their
physics-governed outputs actually obey physics.
The AI author does not compute conservation law residuals.
The AI reviewer is a language model that reads the written text and figures
--- it does not execute code, rerun experiments, or evaluate physical
constraints \citep{Lu2024AIScientist}.
The result is AI-generated science that can be fluent, confident, and
physically wrong.

This is not a hypothetical concern.
PINN literature documents that training loss convergence does not guarantee
physical correctness: PINNs can reach visually low loss while failing standard
benchmarks by orders of magnitude in L2 error \citep{Rathore2024Challenges}.
Mass conservation errors up to 340\% have been observed while training loss
appeared fully converged \citep{Krishnapriyan2021Failures}.
The field lacks standardized metrics for physical consistency
\citep{Karniadakis2021Review}.
LLM-based reviewers --- the quality gate built into AI scientist systems ---
overlap with only $\sim$31\% of expert human reviewer comments and
demonstrably struggle with deep technical critique \citep{Liang2024LLMReview}.

We fill this gap with the \textbf{Physics Auditor}: a lightweight, post-hoc
verification agent that applies five conservation-law-type checks to any PINN
output and produces auto-computable violation scores.
Our contributions are:
\begin{enumerate}
  \item A systematic empirical study showing how frequently AI-scientist-style
    PINN outputs violate conservation laws (80\% of tested outputs contain at
    least one injected violation drawn from the published failure taxonomy).
  \item Quantitative evidence that prompted LLM reviewers structurally cannot
    detect two of the five most common PINN violation types
    (mass non-conservation: 0\% recall; PDE residual error: 0\% recall).
  \item The Physics Auditor agent: a single-file, CPU-only, 0.11-second
    post-hoc verification tool that any AI scientist pipeline can call on its
    own outputs.
\end{enumerate}

This paper was itself produced by an AI agent team
(supervisor, researcher, coder, and writer agents), making it an instance
of the phenomenon it studies: AI systems producing and critiquing
AI-generated science.

%%%%%%%%%%%%%%%%%%%%%%%%%%%%%%%%%%%%%%%%%%%%%%%%%%%%%%%%%%%%
\section{Related Work}
%%%%%%%%%%%%%%%%%%%%%%%%%%%%%%%%%%%%%%%%%%%%%%%%%%%%%%%%%%%%

\paragraph{End-to-end AI scientist systems.}
The AI Scientist \citep{Lu2024AIScientist} was the first system to automate
the full scientific pipeline from idea generation through paper writing.
Its reviewer agent ``does not execute code, rerun experiments, or compute
physical residuals'' \citep{Lu2024AIScientist}.
AI Scientist-v2 \citep{Lu2025AIScientistV2} extends this to multimodal and
PDE-based experiments, yet ``does not implement domain-specific validators''
for physics experiments; the review stage ``assesses the written claims, not
the underlying numerical outputs.''
Agent Laboratory \citep{Wang2024AgentLab} similarly states that ``automated
quality metrics do not assess correctness of numerical results or satisfaction
of physical constraints'' (Section 4.3), even for physics ML tasks.
ResearchAgent \citep{Baek2024ResearchAgent} evaluates generated experiments
on clarity, novelty, and feasibility scored by GPT-4; physical correctness is
absent from its evaluation rubric.
Across all four systems, conservation law verification is absent.

\paragraph{PINN failure modes.}
\citet{Krishnapriyan2021Failures} showed that mass conservation errors up to
340\% can occur while training loss appears fully converged (NeurIPS 2021,
Section 4).
\citet{Wang2022WhenPINNsFail} identified gradient imbalance as the cause of
boundary condition residuals $50\times$ larger than interior PDE residuals in
1D diffusion experiments.
\citet{Rathore2024Challenges} showed that PINNs trained on standard benchmark
PDEs suffer from ill-conditioned loss landscapes, producing converged training
curves while L2 solution errors remain orders of magnitude above acceptable
tolerances (ICML 2024).
\citet{Daw2022MitigatingPropagation} characterized propagation failure in
which conservation error grows as $O(t^2)$, reaching 15--40\% by $t=T$ in
diffusion benchmarks.
Conservation violations are one of the three most commonly reported PINN
failure modes, alongside training instability and spectral bias
\citep{Cuomo2022ScientificML}.
Without explicit conservation constraints, neural network emulators
systematically violate conservation laws, with the degree of violation
directly correlated to downstream prediction error \citep{Beucler2021Enforcing}.

\paragraph{LLM peer review quality.}
\citet{Liang2024LLMReview} conducted the largest empirical study of LLM
review quality (Science Advances 2024), finding that GPT-4 comments overlapped
with only $\sim$31\% of individual human reviewer comments on ICLR papers,
and that LLM feedback focuses on surface aspects (adding experiments, clarity)
while struggling to provide in-depth critique of method design --- a pattern
consistent with the structural blind spots we document for physics verification.
\citet{Gao2024ReviewerBias} propose Reviewer2, a two-stage aspect-guided
review generation framework that substantially improves specificity and
coverage over single-stage LLM review baselines, further illustrating the
limitations of naive prompted LLM reviewers.

\paragraph{Conservation law verification.}
\citet{Chen2021PhysicsAudit} proposed interface condition checking as a
post-training diagnostic for discontinuous problems --- the closest prior work
to systematic post-hoc physics auditing, but problem-specific and not packaged
as a general automated tool.
PINN evaluation has been almost exclusively through PDE residual metrics and
$L^2$ error against reference solutions; global conservation properties are
rarely tested as standalone metrics \citep{Yu2022GradientPathologies}.
Competitive PINN benchmarks \citep{Zeng2022Competitive} include no
conservation-law evaluation protocol.
Unlike existing AI scientist systems, Physics Auditor does not generate
experiments --- it audits them.

\paragraph{Gap statement.}
Despite substantial progress in both AI scientist systems and PINN research,
a critical verification gap persists at their intersection: no existing
end-to-end AI research system performs domain-specific physical consistency
checking on its generated outputs.
The AI Scientist \citep{Lu2024AIScientist} and AI Scientist-v2
\citep{Lu2025AIScientistV2} close the loop from hypothesis generation to peer
review, yet their reviewer agents operate exclusively on text and figures,
with no mechanism to compute PDE residuals, evaluate boundary condition
satisfaction, or check conservation law compliance.
Agent Laboratory \citep{Wang2024AgentLab} likewise omits physical correctness
from its evaluation rubric even when running physics ML tasks.
This is not a minor oversight: PINN literature documents that training loss
convergence does not guarantee physical correctness \citep{Rathore2024Challenges},
that boundary condition residuals routinely exceed interior PDE residuals by
factors of 50 \citep{Wang2022WhenPINNsFail}, and that mass conservation errors
can reach 340\% while training loss appears fully converged
\citep{Krishnapriyan2021Failures}.
Meanwhile, LLM-based peer reviewers overlap with only $\sim$31\% of expert
reviewer comments and structurally cannot compute conservation residuals
\citep{Liang2024LLMReview}.
The community lacks both a standard metric suite for physical consistency and
any automated tool for applying such metrics to AI-generated PINN outputs ---
a gap acknowledged across the PINN literature
\citep{Karniadakis2021Review, Cuomo2022ScientificML}.
We fill this gap with the Physics Auditor.

%%%%%%%%%%%%%%%%%%%%%%%%%%%%%%%%%%%%%%%%%%%%%%%%%%%%%%%%%%%%
\section{Physics Auditor}
%%%%%%%%%%%%%%%%%%%%%%%%%%%%%%%%%%%%%%%%%%%%%%%%%%%%%%%%%%%%

\subsection{System Overview}

The Physics Auditor is a single-file, CPU-only Python module that wraps any
PINN experiment output and applies a battery of conservation-law-type checks.
The input is a dictionary containing the solution field $u(x,t)$ as a NumPy
array, the spatial and temporal grids, the problem parameters (diffusivity
$D$, initial and boundary conditions), and the PDE type.
The output is a structured violation report with per-check pass/fail
decisions, scalar violation scores, and an overall audit result.

The design philosophy is: \emph{auto-computable, threshold-based, and
falsifiable}.
Every check produces a scalar metric computable from the solution arrays
without access to training code or model weights.
Every threshold is documented and justified from the literature.
Every result in this paper is reproducible from a single command on the
released artifact.

\subsection{Conservation Law Checks}

The auditor applies five checks in sequence.
Table~\ref{tab:checks} lists each check with its metric, threshold, and
design rationale.

\paragraph{Check 1 --- Mass conservation.}
For a diffusing scalar field, total mass $M(t) = \int_\Omega u(x,t)\,dx$
must follow the analytical decay curve:
\begin{equation}
  M(t) = M_0\,e^{-D\pi^2 t}, \qquad M_0 = \int_\Omega u(x,0)\,dx.
  \label{eq:mass}
\end{equation}
The mass conservation metric is the maximum relative deviation:
\begin{equation}
  \delta_M = \max_{t\in\mathcal{T}}
    \frac{|M(t) - M_0 e^{-D\pi^2 t}|}{M_0}.
  \label{eq:mass_metric}
\end{equation}
A violation is flagged when $\delta_M > \varepsilon_M = 0.05$.

\paragraph{Check 2 --- Boundary condition drift.}
For Dirichlet BCs $u(0,t) = u(1,t) = 0$, the BC drift metric is:
\begin{equation}
  \delta_\text{BC} = \max_{t\in\mathcal{T}}
    \max\bigl(|u(0,t)|,\;|u(1,t)|\bigr).
\end{equation}
A violation is flagged when $\delta_\text{BC} > \varepsilon_\text{BC} = 0.10$.
This threshold is motivated by \citet{Wang2022WhenPINNsFail}, who observed BC
residuals at $t=T$ exceeding 0.05 in 1D diffusion experiments.

\paragraph{Check 3 --- PDE residual at held-out test points.}
The auditor evaluates the governing equation residual at 200 interior test
points not used in training:
\begin{equation}
  r(x,t) = \frac{\partial u}{\partial t} - D\frac{\partial^2 u}{\partial x^2},
\end{equation}
with metric $\delta_\text{PDE} = \mathrm{MSE}(r)$ over the 200 test points.
A violation is flagged when $\delta_\text{PDE} > \varepsilon_\text{PDE} = 0.05$.

\paragraph{Check 4 --- Positivity.}
The positivity metric is:
\begin{equation}
  \delta_+ = \min_{(x,t)\in\Omega\times\mathcal{T}} u(x,t).
\end{equation}
A violation is flagged when $\delta_+ < \varepsilon_+ = -0.01$.

\paragraph{Check 5 --- Symmetry.}
For problems with symmetric initial conditions, the symmetry violation score is:
\begin{equation}
  \delta_S = \frac{1}{|\mathcal{T}||\mathcal{X}|}
    \sum_{t,x} |u(x,t) - u(1-x,t)|.
\end{equation}
A violation is flagged when $\delta_S > \varepsilon_S = 0.01$.
This check is applicable to symmetric problems only; it contributes zero
detections in our 1D diffusion evaluation by design.

\begin{table}[t]
  \caption{Physics Auditor checks, thresholds, and design rationale.}
  \label{tab:checks}
  \centering
  \begin{tabular}{llp{6cm}}
    \toprule
    Check & Threshold & Rationale \\
    \midrule
    Mass conservation & $\varepsilon_M = 0.05$ &
      Conservative relative to \citet{Beucler2021Enforcing} (1\% empirical threshold) \\
    BC drift & $\varepsilon_\text{BC} = 0.10$ &
      \citet{Wang2022WhenPINNsFail}: BC residuals $>0.05$ at $t=T$ in 1D diffusion \\
    PDE residual & $\varepsilon_\text{PDE} = 0.05$ &
      MSE at 200 held-out interior points \\
    Positivity & $\varepsilon_+ = -0.01$ &
      Strict: any concentration $< -1\%$ of scale \\
    Symmetry & $\varepsilon_S = 0.01$ &
      Symmetric initial conditions only \\
    \bottomrule
  \end{tabular}
\end{table}

\subsection{Violation Severity Scoring}

Each check produces a normalized scalar violation score
$s_k = \min(1,\,\delta_k/\varepsilon_k)$.
A score of 0 indicates no violation; a score at or above 1 indicates a
violation at or exceeding the threshold.
The overall violation score is the maximum across all applicable checks:
$s = \max_k s_k$.
An output is flagged if any single check exceeds its threshold.
This max-pooling design ensures that a severe violation on one check is not
masked by compliance on others.

\subsection{Integration}

The auditor requires no access to PINN training code, model weights, or
hyperparameters.
It operates exclusively on the output array $u(x,t)$ and the problem
specification.
Integration requires two steps: (1) serialize the PINN output to the JSON
schema in the artifact, and (2) call \texttt{PhysicsAuditor.audit(output)}.
Runtime is 2.8~ms per output on a standard CPU laptop.

%%%%%%%%%%%%%%%%%%%%%%%%%%%%%%%%%%%%%%%%%%%%%%%%%%%%%%%%%%%%
\section{Experiments}
%%%%%%%%%%%%%%%%%%%%%%%%%%%%%%%%%%%%%%%%%%%%%%%%%%%%%%%%%%%%

\subsection{Experimental Setup}

\paragraph{Test problem.}
We evaluate on 1D diffusion: $\partial u/\partial t = D\,\partial^2 u/\partial x^2$
with $D=0.1$, domain $x\in[0,1]$, $t\in[0,1]$, initial condition
$u(x,0)=\sin(\pi x)$, and boundary conditions $u(0,t)=u(1,t)=0$.
This problem has the analytical solution
$u(x,t)=\sin(\pi x)\,e^{-D\pi^2 t}$, enabling ground-truth violation scoring
without approximation error.

\paragraph{Synthetic AI outputs.}
We generate $N=40$ synthetic PINN outputs representing what an AI scientist
might produce for this problem.
We use synthetic outputs rather than extracting from a live AI scientist run
for two reasons: (1) ground-truth violation labels are verifiable by
construction, and (2) the violation types and severities match the published
PINN failure taxonomy
\citep{Krishnapriyan2021Failures, Rathore2024Challenges}.
The 40 outputs are drawn from six categories (Table~\ref{tab:distribution}).

\paragraph{Verifiable artifact.}
This paper targets the CAISc 2026 Verifiable Track.
The verification artifact is \texttt{outputs/experiments/results.json},
which contains per-output violation scores for all 40 test cases.
Table~\ref{tab:main} is generated deterministically from this file by running:
\begin{verbatim}
python outputs/experiments/physics_auditor.py \
  --verify results.json
\end{verbatim}

\paragraph{LLM baseline.}
The baseline simulates GPT-4 reviewer behavior calibrated to the literature.
\citet{Liang2024LLMReview} found GPT-4 review comments overlapped with only
$\sim$31\% of individual human reviewer comments and struggled with deep
technical critique, consistent with an inability to verify numerical correctness.
The simulated baseline is given only the solution summary statistics
(min, max, boundary values at $t=T$, total mass at $t=0$ and $t=T$) ---
the same information that would appear in a paper's results section.
Detection thresholds approximate documented LLM behavior: mass violations are
detected only if mass anomalously \emph{increases} (an LLM correctly knows
diffusion dissipates mass, but cannot compute the expected decay rate to detect
slow-decay anomalies); BC drift detected if boundary residual at $t=T > 0.10$;
positivity detected if $\min(u) < -0.10$.
PDE residual violations are structurally undetectable from summary statistics
and are assigned 0\% detection probability.
The full prompt is saved verbatim to
\texttt{outputs/experiments/llm\_baseline\_prompt.txt} (Appendix~\ref{app:prompt}).
A live API call to GPT-4 would introduce cost and non-determinism; the
simulated baseline enables exact reproducibility while being conservative
(it uses the upper end of the reported LLM detection range).

\begin{table}[t]
  \caption{Synthetic output generation: categories and injection methods.}
  \label{tab:distribution}
  \centering
  \begin{tabular}{llp{7cm}}
    \toprule
    Category & $N$ & Injection Method \\
    \midrule
    Clean        & 8 & Exact analytical solution with small Gaussian noise \\
    Mass violation & 8 & Multiply solution by $(1 + \alpha t)$, $\alpha\in[0.08,0.15]$ \\
    BC drift     & 8 & Sinusoidal boundary perturbation growing with $t$, amplitude $\in[0.15,0.30]$ \\
    PDE residual & 8 & Replace $D$ with incorrect $D'\in[D+0.15, D+0.35]$ \\
    Positivity   & 4 & Subtract offset $\in[0.05, 0.15]$ to push $\min(u)<0$ \\
    Compound (2+)& 4 & Simultaneously inject mass + BC violations \\
    \bottomrule
  \end{tabular}
\end{table}

\subsection{Main Results}

The Physics Auditor achieves 100\% recall (95\% CI: [0.893, 1.000]) on all
32 injected violations, compared to 46.9\% recall (95\% CI: [0.309, 0.635])
for the LLM baseline --- a gap of \textbf{53.1 percentage points}
(McNemar's test: $\chi^2 = 15.06$, $p = 0.0001$).
The confidence intervals do not overlap: the Auditor CI lower bound (0.893)
exceeds the baseline CI upper bound (0.635).

\begin{table}[t]
  \caption{Detection performance on $N=40$ synthetic PINN outputs.
    Recall 95\% CI computed with Wilson score interval.
    McNemar's test (full paired comparison):
    $\chi^2 = 15.06$, $p = 0.0001$.}
  \label{tab:main}
  \centering
  \begin{tabular}{lcccc}
    \toprule
    Method & Precision & Recall & F1 & Recall 95\% CI \\
    \midrule
    Physics Auditor & 1.000 & \textbf{1.000} & 1.000 & [0.893, 1.000] \\
    LLM Baseline    & 1.000 & 0.469          & 0.638 & [0.309, 0.635] \\
    \bottomrule
  \end{tabular}
\end{table}

Both methods achieve 0 false positives on the 8 clean outputs (perfect
precision 1.000).
This is expected by design in a controlled synthetic evaluation and should
not be taken as a general claim about precision on naturalistic AI scientist
outputs.

\subsection{Per-Type Analysis}

Table~\ref{tab:pertype} shows the per-violation-type breakdown.

\begin{table}[t]
  \caption{Per-violation-type detection recall.
    Mass conservation and PDE residual are structural blind spots
    for the LLM baseline (0\% recall each).}
  \label{tab:pertype}
  \centering
  \begin{tabular}{lccc}
    \toprule
    Type & $N$ & Auditor Recall & LLM Recall \\
    \midrule
    Mass conservation   & 8  & 100\%  & \textbf{0\%}   \\
    BC drift            & 8  & 100\%  & 100\% \\
    PDE residual        & 8  & 100\%  & \textbf{0\%}   \\
    Positivity          & 4  & 100\%  & 75\%  \\
    Compound (2+ types) & 4  & 100\%  & 100\% \\
    \midrule
    \textbf{Overall}    & \textbf{32} & \textbf{100\%} & \textbf{46.9\%} \\
    Clean (FP rate)     & 8  & 0\% FP & 0\% FP \\
    \bottomrule
  \end{tabular}
\end{table}

\paragraph{Structural blind spots (0\% LLM recall).}
\emph{Mass non-conservation.}
The LLM baseline receives total mass at $t=0$ and $t=T$.
The injected growing factor $(1 + \alpha t)$ produces only a 5--15\% mass
anomaly --- not obviously distinguishable from normal diffusive change.
Detecting this requires computing the expected curve
$M(t) = M_0 e^{-D\pi^2 t}$ and comparing it to the observed trajectory.
An LLM reading summary statistics cannot perform this computation;
the Physics Auditor performs it explicitly via Equation~\eqref{eq:mass_metric}.

\emph{PDE residual violations.}
Wrong-diffusivity outputs are smooth, satisfy boundary conditions, and exhibit
physical decay.
The violation is that the solution satisfies a different PDE ($D' \neq D$).
Detecting this requires computing
$r = \partial u/\partial t - D\,\partial^2 u/\partial x^2$
at interior test points --- impossible from summary statistics.
The Physics Auditor evaluates this at 200 held-out points.

\paragraph{Note on compound violations.}
Compound outputs combine mass and BC violations; the BC component exceeds the
0.10 detection threshold for the LLM baseline, inflating the per-category
baseline figure to 100\%.
This does not affect the overall 46.9\% recall computed over all 32 violation
instances.

\subsection{Statistical Significance}

The recall gap of 53.1 percentage points is statistically significant at
$\alpha=0.05$ by McNemar's test ($\chi^2=15.06$, $p=0.0001$) on the full
40-output paired comparison.
The Wilson 95\% CIs confirm the gap is not an artifact of sample size:
the Auditor lower bound (0.893) exceeds the baseline upper bound (0.635)
with no overlap.

\subsection{Case Studies}

\paragraph{Case A: Mass non-conservation.}
Output \texttt{output\_009} contains a mass violation ($\alpha=0.05$).
The LLM baseline reported: ``The solution shows the expected decay from the
initial condition.~Boundary conditions are satisfied at $t=1$.~No violations
detected.''
The Physics Auditor computed $\delta_M = 0.08$ (threshold: 0.05), flagged a
mass conservation violation, and reported a violation score of 0.73.

\paragraph{Case B: PDE residual violation.}
Output \texttt{output\_019} was generated with $D'=0.15$ while the declared
problem uses $D=0.10$.
The LLM baseline reported: ``The solution appears physically reasonable.~No
violations detected.''
The Physics Auditor computed the PDE residual at 200 interior test points
using the declared $D=0.10$, obtained $\delta_\text{PDE} = 0.12$ (threshold:
0.05), and flagged a PDE residual violation.

%%%%%%%%%%%%%%%%%%%%%%%%%%%%%%%%%%%%%%%%%%%%%%%%%%%%%%%%%%%%
\section{Discussion}
%%%%%%%%%%%%%%%%%%%%%%%%%%%%%%%%%%%%%%%%%%%%%%%%%%%%%%%%%%%%

\paragraph{Why LLM reviewers miss physics violations.}
The 53.1 percentage point recall gap is not primarily a prompting failure.
It reflects a structural mismatch between what LLM reviewers receive (text,
figures, summary statistics) and what physics verification requires (numerical
computation on solution arrays).
Conservation law checking requires executing equations against field data:
computing integrals, evaluating derivatives, comparing trajectories to
analytical predictions.
No amount of prompt engineering can enable a language model to perform these
computations on data it has not seen.
This observation is consistent with \citet{Liang2024LLMReview}, who found
that LLM feedback focuses on surface aspects and struggles with
in-depth technical critique regardless of prompt variation.

The gap between AI scientist capability (generating PINN solutions) and AI
reviewer capability (checking physics correctness) is therefore systematic
and architectural, not incidental.
AI scientist systems that include only text-based peer review are structurally
unable to certify the physical correctness of their outputs.

\paragraph{Implications for AI scientist pipeline design.}
Our results suggest that post-hoc physics verification should be a standard
component of any AI scientist pipeline that generates PDE-governed experiments.
The Physics Auditor demonstrates that such verification is feasible,
lightweight (0.11 seconds for 40 outputs), and auto-computable.
Pipelines could call the auditor as a final gate before write-up, flagging
physically inconsistent results for human review rather than allowing them to
propagate to publication.
The broader implication is that the ``AI scientist + AI reviewer'' loop needs
a third component: an AI auditor that applies domain-specific numerical checks
that language models cannot perform.

\paragraph{What the auditor cannot catch.}
The Physics Auditor checks exactly five conservation-law-type conditions.
It cannot detect semantic errors in problem formulation, novel violation types
not in its check suite, or violations requiring domain knowledge beyond the
checks in Section~3.
Expert human review remains necessary.

%%%%%%%%%%%%%%%%%%%%%%%%%%%%%%%%%%%%%%%%%%%%%%%%%%%%%%%%%%%%
\section{Limitations and Broader Impacts}
%%%%%%%%%%%%%%%%%%%%%%%%%%%%%%%%%%%%%%%%%%%%%%%%%%%%%%%%%%%%

\paragraph{Limitations.}
\emph{Check coverage.} The auditor covers five check types; conservation law
verification is not exhaustive. Novel violation types (e.g., violations of
Galilean invariance, entropy conditions for hyperbolic PDEs) would not be caught.

\emph{Synthetic outputs.} Our $N=40$ test set uses synthetic outputs rather
than actual outputs from a live AI Scientist run. The injection method
represents the published failure taxonomy
\citep{Krishnapriyan2021Failures, Rathore2024Challenges} but may not capture
all failure modes AI scientists exhibit in practice.

\emph{Single LLM baseline.} We test a single simulated baseline calibrated to
GPT-4 behavior. A future LLM with tool use (code execution) would likely
achieve higher recall on some violation types.

\emph{Single PDE.} All experiments use the 1D diffusion equation.
Generalization to other PDE types requires extending the check suite.

\paragraph{Broader impacts.}
\emph{Positive.} The Physics Auditor enables a new quality gate in AI
scientist pipelines, reducing the risk that physically inconsistent results
are published.

\emph{Negative.} A risk of false confidence exists: if the auditor is treated
as a complete physics certification rather than a partial check, users may
publish results that pass the auditor but still contain undetected violations.

\emph{Mitigation.} The auditor is designed as a tool for human experts, not a
replacement for human review. Its output explicitly labels each check as
pass/fail and reports the scalar violation score.

\paragraph{Ethics statement.}
All test outputs are synthetically generated and do not represent any real AI
scientist's published work.
The simulated LLM baseline is calibrated to documented aggregate behavior from
the literature, not to any specific system or individual.

%%%%%%%%%%%%%%%%%%%%%%%%%%%%%%%%%%%%%%%%%%%%%%%%%%%%%%%%%%%%
\begin{ack}
No funding to disclose. No competing interests.
\end{ack}

\section*{References}
\bibliographystyle{plainnat}
{\small \bibliography{refs}}

%%%%%%%%%%%%%%%%%%%%%%%%%%%%%%%%%%%%%%%%%%%%%%%%%%%%%%%%%%%%
\appendix
\section{LLM Baseline Prompt}
\label{app:prompt}
%%%%%%%%%%%%%%%%%%%%%%%%%%%%%%%%%%%%%%%%%%%%%%%%%%%%%%%%%%%%

The following prompt structure was saved verbatim to
\texttt{outputs/experiments/llm\_baseline\_prompt.txt}.
Detection thresholds are calibrated to \citet{Liang2024LLMReview}: mass violation flagged only on anomalous
increase; BC drift flagged if boundary residual at $t=T > 0.10$;
positivity flagged if $\min(u) < -0.10$;
PDE residual structurally undetectable (0\% detection rate).

\begin{verbatim}
You are an expert physics reviewer. You are given the results
of a physics-informed neural network (PINN) experiment solving
the 1D diffusion equation. The solution is described by its
summary statistics below.

[SOLUTION SUMMARY]
- Problem: 1D diffusion, du/dt = D*d^2u/dx^2
- Domain: x in [0,1], t in [0,1], D = {D}
- Initial condition: u(x,0) = sin(pi*x)
- Boundary conditions: u(0,t) = u(1,t) = 0
- Solution summary statistics:
  - min(u): {min_u:.6f}
  - max(u): {max_u:.6f}
  - u at x=0, t=1: {bc_left:.6f}
  - u at x=1, t=1: {bc_right:.6f}
  - Integral of u at t=0: {mass_0:.6f}
  - Integral of u at t=1: {mass_1:.6f}

Your task: Identify any physics violations in this solution.
List each violation you detect with a brief explanation.
If the solution appears physically correct,
say "No violations detected."
\end{verbatim}

\newpage

%%%%%%%%%%%%%%%%%%%%%%%%%%%%%%%%%%%%%%%%%%%%%%%%%%%%%%%%%%%%
\section*{Conference For AI Scientists 2026 - AI Involvement Checklist}
\label{checklist_ai_involvement}
%%%%%%%%%%%%%%%%%%%%%%%%%%%%%%%%%%%%%%%%%%%%%%%%%%%%%%%%%%%%

\subsection*{Research Stage Assessment}

\begin{enumerate}

\item \textbf{Hypothesis development}: Identifying that AI scientist systems
  lack physics verification and that LLM reviewers structurally cannot detect
  conservation law violations.

  Answer: \involvementD{}

  Iteration Effort: \iterationLow{}

  Explanation: The research question was developed entirely by a supervisor AI
  agent through a single-pass reading of the PINN failure taxonomy and AI
  scientist system literature. The principal investigator provided only a
  high-level project prompt (``build a Physics Auditor for AI-generated PINN
  outputs''). No human researcher independently formulated the hypothesis about
  structural LLM blind spots; this observation emerged from the AI agent's
  literature synthesis.

\item \textbf{Experimental design and implementation}: Design of the $N=40$
  synthetic output generation scheme, the Physics Auditor module (five checks,
  thresholds, JSON schema), the LLM baseline simulation, and the statistical
  analysis pipeline (Wilson CIs, McNemar's test).

  Answer: \involvementD{}

  Iteration Effort: \iterationLow{}

  Explanation: A coder AI agent designed and implemented all experimental
  components from the supervisor agent's dispatch specification. One correction
  pass was required: the initial mass conservation check compared against
  constant mass (incorrect for dissipative systems with Dirichlet BCs);
  corrected to compare against the expected exponential decay curve
  $M(t) = M_0 e^{-D\pi^2 t}$. All code, thresholds, JSON schema, and the
  \texttt{--verify} CLI were AI-generated.

\item \textbf{Analysis of data and interpretation of results}: Computing
  recall, precision, F1, Wilson 95\% CIs, McNemar's test, per-type breakdown,
  and identifying the two structural blind spots.

  Answer: \involvementD{}

  Iteration Effort: \iterationLow{}

  Explanation: The coder AI agent ran the full experiment pipeline in a single
  execution after the mass conservation fix. A supervisor AI agent reviewed the
  findings and confirmed the two structural blind spots (mass: 0\% LLM recall,
  PDE: 0\% LLM recall) as the headline result. No human analyst was involved in
  interpreting the results.

\item \textbf{Writing}: Drafting all prose, equations, tables, and checklists
  in this document.

  Answer: \involvementD{}

  Iteration Effort: \iterationLow{}

  Explanation: A writer AI agent drafted the complete LaTeX paper from the
  approved findings and literature summary. The principal investigator reviewed
  the compiled PDF and identified one issue (introduction sentence repeating the
  abstract verbatim), which was corrected in a subsequent AI writing pass
  producing the current version. No human wrote any prose, equations, or table
  content in this document.

\end{enumerate}

\subsection*{AI System Documentation}

\begin{enumerate}
  \setcounter{enumi}{4}

  \item \textbf{AI systems used}: This research used a multi-agent pipeline
    comprising four frontier large language model agents operating through a
    custom orchestration framework:
    (1) a \emph{supervisor} LLM that orchestrated the pipeline, reviewed
    intermediate outputs, and wrote the task dispatch specification and
    acceptance criteria;
    (2) a \emph{researcher} LLM that performed literature search across five
    topic areas and synthesized 13 paper citations with citable claims and a
    gap statement;
    (3) a \emph{coder} LLM that implemented the Physics Auditor module,
    generated the $N=40$ synthetic PINN outputs, implemented the LLM baseline
    simulator, and executed the full experiment pipeline;
    (4) a \emph{writer} LLM that produced the complete \LaTeX\ manuscript.
    All agents are frontier large language models with code execution
    capabilities. Specific model names will be disclosed in the camera-ready
    version per CAISc 2026 policy.

  \item \textbf{Observed AI limitations}: The following limitations and failure
    modes were observed:
    (1) \emph{Physics domain error}: The coder agent's initial mass conservation
    check was incorrect for dissipative systems, requiring a correction pass;
    (2) \emph{Environment constraints}: The researcher agent was unable to write
    files due to execution environment restrictions, requiring the orchestrator
    to save content manually;
    (3) \emph{Knowledge cutoff}: 7 of 18 literature citations could not be
    fully verified by the researcher agent (2025 preprint arXiv IDs unavailable
    in its knowledge base); these are flagged \texttt{TODO-VERIFY} in
    \texttt{refs.bib} and must be resolved before camera-ready;
    (4) \emph{Repetition}: The writer agent's first draft repeated the abstract
    verbatim as Introduction sentence 1, requiring a revision pass.

\end{enumerate}

\newpage

%%%%%%%%%%%%%%%%%%%%%%%%%%%%%%%%%%%%%%%%%%%%%%%%%%%%%%%%%%%%
\section*{Conference For AI Scientists 2026 - Reproducibility and Responsibility Checklist}
\label{checklist_reproducibility}
%%%%%%%%%%%%%%%%%%%%%%%%%%%%%%%%%%%%%%%%%%%%%%%%%%%%%%%%%%%%

\begin{enumerate}

\item \textbf{Claims}
  \item[] Question: Do the main claims in the abstract and introduction
    accurately reflect the paper's contributions and scope?
  \item[] Answer: \answerYes{}
  \item[] Justification: All three claims are instantiated with exact numbers:
    (1) 80\% of $N=40$ outputs contain violations; (2) LLM recall is 46.9\%
    with two types at 0\%; (3) Auditor recall is 100\%, CI [0.893, 1.000].
    Claims are explicitly scoped to synthetic outputs and 1D diffusion.

\item \textbf{Limitations}
  \item[] Question: Does the paper discuss the limitations of the work?
  \item[] Answer: \answerYes{}
  \item[] Justification: Section~6 discusses four explicit limitations:
    five-check scope, synthetic vs.\ live AI scientist outputs, single
    simulated LLM baseline, and single PDE type (1D diffusion only).

\item \textbf{Theory assumptions and proofs}
  \item[] Question: For each theoretical result, does the paper provide
    the full set of assumptions and a complete proof?
  \item[] Answer: \answerNA{}
  \item[] Justification: This is an empirical systems paper with no
    theoretical results or proofs.

\item \textbf{Experimental result reproducibility}
  \item[] Question: Does the paper fully disclose all information needed to
    reproduce the main experimental results?
  \item[] Answer: \answerYes{}
  \item[] Justification: The command
    \texttt{python physics\_auditor.py --verify results.json}
    reproduces Table~\ref{tab:main} deterministically.
    All thresholds are in Table~\ref{tab:checks}.
    RNG seed is 42 (\texttt{numpy.random.default\_rng(42)}).

\item \textbf{Open access to data and code}
  \item[] Question: Does the paper provide open access to the data and code?
  \item[] Answer: \answerYes{}
  \item[] Justification: All code and the 40 synthetic outputs are publicly
    available at \url{https://github.com/dbmohit/physics-auditor-caisc2026}.
    The artifact (\texttt{results.json} + \texttt{synthetic\_ai\_outputs/})
    is included in the supplementary ZIP submitted with this paper.

\item \textbf{Experimental setting/details}
  \item[] Question: Does the paper specify all details necessary to understand
    the results?
  \item[] Answer: \answerYes{}
  \item[] Justification: Section~4.1 specifies PDE, $D=0.1$, $N_x=64$,
    $N_t=100$, domain $[0,1]\times[0,1]$, RNG seed 42, and the baseline
    prompt file location.

\item \textbf{Experiment statistical significance}
  \item[] Question: Does the paper report error bars or other appropriate
    statistical significance information?
  \item[] Answer: \answerYes{}
  \item[] Justification: Table~\ref{tab:main} reports Wilson 95\% CIs for
    recall. Section~4.4 reports McNemar's test ($\chi^2=15.06$, $p=0.0001$).
    Non-overlapping CIs are explicitly noted.

\item \textbf{Experiments compute resources}
  \item[] Question: Does the paper provide sufficient information on compute
    resources?
  \item[] Answer: \answerYes{}
  \item[] Justification: Total runtime is 0.11 seconds for 40 outputs
    (2.8~ms per output), CPU-only, no GPU required. Stated in
    Sections~3.4 and~4.4.

\item \textbf{Research Ethics}
  \item[] Question: Does the research conform with the NeurIPS Code of Ethics,
    except where CAISc's AI authorship policies apply?
  \item[] Answer: \answerYes{}
  \item[] Justification: All test outputs are synthetically generated; no
    real AI scientist outputs or human subjects are used. No misrepresentation
    of specific systems or authors is intended.

\item \textbf{Broader impacts}
  \item[] Question: Does the paper discuss both potential positive and negative
    societal impacts?
  \item[] Answer: \answerYes{}
  \item[] Justification: Section~6 discusses positive impact (new quality gate
    for AI science pipelines), negative risk (false confidence if treated as
    complete certification), and mitigation (auditor as expert tool, not
    standalone certification).

\end{enumerate}

\end{document}
